% GENERATED FILE. Edit canonical lessons, then run npm run generate:books.
% canonical-source-hash: fnv1a64:fa949d984dcca324
% canonical-lessons: FR-C392-especes, FR-C392-cheque, FR-C392-chequier, FR-C392-virement, FR-C392-retirer

\chapter{À la banque}
\label{ch:fr-a-la-banque}
\hlchaptermodality{392}

\begin{chapteropening}
\textbf{By the end of this chapter:} \emph{I can handle money at the bank: cash, a cheque, a chequebook, a transfer and taking money out.}

Complete the last lesson of chapter 392: I can handle money at the bank: cash, a cheque, a chequebook, a transfer and taking money out.
\end{chapteropening}

\section[les espèces]{les espèces — the cash}

\label{lesson:FR-C392-especes}



\begin{quote}
Before the new one: say the French for the investment, then the French for the rate.
\end{quote}

\subsection*{You'll want to know: les espèces}
\textbf{les espèces} — \textquotedblleft{}the cash\textquotedblright{}.

\begin{quote}
Ce petit café n'accepte que les espèces.
\end{quote}

\begin{quote}
\emph{This small café only takes cash.}
\end{quote}

\emph{Payer en espèces} — to pay cash. Québec: \emph{payer comptant}. In the singular, \emph{une espèce} is a species or a kind.

\begin{cousinweb}
\emph{espèces} is the plural of \emph{espèce} \textquotedblleft{}kind\textquotedblright{}, from Latin \emph{speciēs} \textquotedblleft{}appearance, kind\textquotedblright{} (Lewis \& Short), from \emph{specere} \textquotedblleft{}to look\textquotedblright{}: money in a visible form.
\end{cousinweb}

\begin{grammarlens}[title={What you've built}]
One more word for telling, explaining and describing at B1.
\end{grammarlens}

\subsection*{Your turn}
\begin{itemize}
\raggedright
  \item \emph{Say it:} \emph{les espèces}
  \item \emph{Say it:} \emph{les espèces}, once more
  \item \emph{Say it:} a sentence of your own with \emph{les espèces}
  \item \emph{Recall:} say the French for the investment, then the French for the rate, then say \emph{les espèces} again
\end{itemize}

\begin{tcolorbox}[breakable,colback=teal!4,colframe=teal!35!black,title={Before you move on}]
What does les espèces mean? (\textquotedblleft{}The cash\textquotedblright{}.) Say it once more.
\end{tcolorbox}
\section[le chèque]{le chèque — the cheque}

\label{lesson:FR-C392-cheque}



\begin{quote}
Before the new one: say the French for the rate, then the French for the cash.
\end{quote}

\subsection*{You'll want to know: le chèque}
\textbf{le chèque} — \textquotedblleft{}the cheque\textquotedblright{}.

\begin{quote}
Elle a payé le plombier par chèque.
\end{quote}

\begin{quote}
\emph{She paid the plumber by cheque.}
\end{quote}

\emph{Un chèque sans provision} — a bounced cheque.

\begin{cousinweb}
\emph{chèque} is usually traced to English \emph{cheque}, \emph{check}, itself linked to \emph{check} \textquotedblleft{}a control\textquotedblright{}: a paper that guards against forgery.
\end{cousinweb}

\begin{grammarlens}[title={What you've built}]
One more word for telling, explaining and describing at B1.
\end{grammarlens}

\subsection*{Your turn}
\begin{itemize}
\raggedright
  \item \emph{Say it:} \emph{le chèque}
  \item \emph{Say it:} \emph{le chèque}, once more
  \item \emph{Say it:} a sentence of your own with \emph{le chèque}
  \item \emph{Recall:} say the French for the rate, then the French for the cash, then say \emph{le chèque} again
\end{itemize}

\begin{tcolorbox}[breakable,colback=teal!4,colframe=teal!35!black,title={Before you move on}]
What does le chèque mean? (\textquotedblleft{}The cheque\textquotedblright{}.) Say it once more.
\end{tcolorbox}
\section[le chéquier]{le chéquier — the chequebook}

\label{lesson:FR-C392-chequier}



\begin{quote}
Before the new one: say the French for the cash, then the French for the cheque.
\end{quote}

\subsection*{You'll want to know: le chéquier}
\textbf{le chéquier} — \textquotedblleft{}the chequebook\textquotedblright{}.

\begin{quote}
La banque m'a envoyé un nouveau chéquier.
\end{quote}

\begin{quote}
\emph{The bank sent me a new chequebook.}
\end{quote}

From \emph{chèque}. Fewer and fewer people use one.

\begin{cousinweb}
\emph{chéquier} is built from \emph{chèque} and the suffix \emph{-ier}. \emph{chèque} is usually traced to English \emph{cheque}, \emph{check}, linked to \emph{check} \textquotedblleft{}a control\textquotedblright{}.
\end{cousinweb}

\begin{grammarlens}[title={What you've built}]
One more word for telling, explaining and describing at B1.
\end{grammarlens}

\subsection*{Your turn}
\begin{itemize}
\raggedright
  \item \emph{Say it:} \emph{le chéquier}
  \item \emph{Say it:} \emph{le chéquier}, once more
  \item \emph{Say it:} a sentence of your own with \emph{le chéquier}
  \item \emph{Recall:} say the French for the cash, then the French for the cheque, then say \emph{le chéquier} again
\end{itemize}

\begin{tcolorbox}[breakable,colback=teal!4,colframe=teal!35!black,title={Before you move on}]
What does le chéquier mean? (\textquotedblleft{}The chequebook\textquotedblright{}.) Say it once more.
\end{tcolorbox}
\section[le virement]{le virement — the bank transfer}

\label{lesson:FR-C392-virement}



\begin{quote}
Before the new one: say the French for the cheque, then the French for the chequebook.
\end{quote}

\subsection*{You'll want to know: le virement}
\textbf{le virement} — \textquotedblleft{}the bank transfer\textquotedblright{}.

\begin{quote}
J'ai fait un virement pour payer le loyer.
\end{quote}

\begin{quote}
\emph{I made a transfer to pay the rent.}
\end{quote}

\emph{Un virement bancaire} — a bank transfer.

\begin{cousinweb}
\emph{virement} is built from \emph{virer} \textquotedblleft{}to turn, to transfer\textquotedblright{} and \emph{-ment} (EtymDB). The origin of \emph{virer} is usually traced to a Latin form linked to \emph{vibrāre} \textquotedblleft{}to shake\textquotedblright{}; it is not settled.
\end{cousinweb}

\begin{grammarlens}[title={What you've built}]
One more word for telling, explaining and describing at B1.
\end{grammarlens}

\subsection*{Your turn}
\begin{itemize}
\raggedright
  \item \emph{Say it:} \emph{le virement}
  \item \emph{Say it:} \emph{le virement}, once more
  \item \emph{Say it:} a sentence of your own with \emph{le virement}
  \item \emph{Recall:} say the French for the cheque, then the French for the chequebook, then say \emph{le virement} again
\end{itemize}

\begin{tcolorbox}[breakable,colback=teal!4,colframe=teal!35!black,title={Before you move on}]
What does le virement mean? (\textquotedblleft{}The bank transfer\textquotedblright{}.) Say it once more.
\end{tcolorbox}
\section[retirer]{retirer — to take out, to withdraw}

\label{lesson:FR-C392-retirer}



\begin{quote}
Before the new one: say the French for the chequebook, then the French for the bank transfer.
\end{quote}

\subsection*{You'll want to know: retirer}
\textbf{retirer} — \textquotedblleft{}to take out, to withdraw\textquotedblright{}.

\begin{quote}
J'ai retiré cinquante euros au distributeur.
\end{quote}

\begin{quote}
\emph{I took fifty euros out of the cash machine.}
\end{quote}

Also to take off: \emph{retirer son manteau}.

\begin{cousinweb}
\emph{retirer} is built from \emph{re-} \textquotedblleft{}back\textquotedblright{} and \emph{tirer} \textquotedblleft{}to pull\textquotedblright{} (EtymDB). The origin of \emph{tirer} is not settled.
\end{cousinweb}

\begin{grammarlens}[title={What you've built}]
One more word for telling, explaining and describing at B1.
\end{grammarlens}

\subsection*{Your turn}
\begin{itemize}
\raggedright
  \item \emph{Say it:} \emph{retirer}
  \item \emph{Say it:} \emph{retirer}, once more
  \item \emph{Say it:} a sentence of your own with \emph{retirer}
  \item \emph{Recall:} say the French for the chequebook, then the French for the bank transfer, then say \emph{retirer} again
\end{itemize}

\begin{tcolorbox}[breakable,colback=teal!4,colframe=teal!35!black,title={Before you move on}]
What does retirer mean? (\textquotedblleft{}To take out, to withdraw\textquotedblright{}.) Say it once more.
\end{tcolorbox}
